\chapter{Laminar Non-Premixed (Diffusion) Flames}\label{ch:nonpremixed}

In a non-premixed flame the fuel and oxidiser arrive from opposite sides and meet only at the flame, so there is no single reactants-to-products coordinate and no intrinsic propagation speed: the problem is one of \emph{mixing}, and the flame sits wherever the local mixture is stoichiometric~\cite{lecture9,transcript_l9}. This chapter introduces the mixture-fraction description that replaces the progress variable, the Burke--Schumann limit, the counterflow flamelet model with its scalar dissipation rate, and the qualitatively different way diffusion flames respond to strain.

\section{The mixture fraction}

Because the progress variable presupposes a reactants-to-products path, it cannot describe a diffusion flame; what is needed instead is a variable that uniquely locates a point relative to the flame in terms of how much of its mass came from each feed stream~\cite{lecture9,transcript_l9}. The \textbf{mixture fraction} $Z$ is defined as the fraction of the local mass that originated in the fuel stream,
\begin{equation}
Z=\frac{m_1}{m_1+m_2},\qquad Z=0\ \text{in pure oxidiser},\quad Z=1\ \text{in pure fuel}.
\label{eq:Zdef}
\end{equation}
For one-step chemistry $\nu_F' F+\nu_O' O\rightarrow P$, the oxidiser-to-fuel stoichiometric mass ratio is $s=Y_O^0/Y_F^0\big|_{st}=\nu_O' W_O/(\nu_F' W_F)$ (for $\ch{CH_4}+2\ch{O_2}$, $s=4$)~\cite{lecture9}. Writing the unburnt fuel and oxidiser mass fractions in terms of $Z$ and using $\dd Y_O/(\nu_O' W_O)=\dd Y_F/(\nu_F' W_F)$ shows that the combination $sY_F-Y_O$ is unaffected by reaction, which yields the explicit, reaction-independent definition
\begin{equation}
Z=\frac{sY_F-Y_O+Y_O^0}{sY_F^0+Y_O^0}.
\label{eq:Z-twostream}
\end{equation}
The decisive property is that $Z$ obeys a transport equation with \emph{no chemical source term},
\begin{equation}
\pder{}{t}(\rho Z)+\divv(\rho Z\vv v)=\divv(\rho D\,\grad Z),
\label{eq:Z-transport}
\end{equation}
so $Z$ is governed purely by convection and diffusion (assuming a single diffusion coefficient for all species)~\cite{lecture9}. This is precisely why non-premixed flames are called \textbf{diffusion flames}: their structure is set by mixing, whereas in premixed flames convection (the flame speed) plays the leading role~\cite{transcript_l9}.

\section{The stoichiometric mixture fraction and the flame location}

The flame sits at the iso-surface where the mixture is exactly stoichiometric, $Z=\Zst$, given by
\begin{equation}
\Zst=\frac{1}{1+sY_F^0/Y_O^0}=\frac{1}{1+\nu_O' W_O\,Y_F^0/(\nu_F' W_F\,Y_O^0)} .
\label{eq:Zst}
\end{equation}
For methane burning in air this is small, $\Zst\approx0.054$, meaning the flame sits close to the oxidiser (air) side, because comparatively little fuel is needed to consume a given mass of the air stream~\cite{transcript_l9}. Two manipulations shift it~\cite{transcript_l9}: \emph{diluting the oxidiser} (replacing pure \ch{O_2} with air, i.e.\ adding \ch{N_2}) moves the flame toward the \emph{fuel} side ($\Zst$ increases), because more oxidiser-stream mass is now needed to supply the same \ch{O_2}; and \emph{changing the fuel} to one with a higher mass-based stoichiometric air requirement likewise pushes $\Zst$ toward the fuel side. In a counterflow geometry the flame accordingly sits on the oxidiser side of the stagnation plane~\cite{transcript_l9}.

\paragraph{Bilger's element-based mixture fraction.}
The two-stream definition of Eq.~\eqref{eq:Z-twostream} is awkward in measurements and breaks down once reaction redistributes the species unless special conditions hold. The robust alternative is \textbf{Bilger's mixture fraction}, built from \emph{element} mass fractions~\cite{lecture9,poinsot2011theoretical}. Since chemical reactions only rearrange elements between species and never change the total mass of carbon, hydrogen or oxygen, element mass fractions are conserved; combining them in the particular linear combination $\beta$ whose reaction-rate contributions cancel, and normalising so that $\beta$ runs from $0$ in the oxidiser to $1$ in the fuel, gives a mixture fraction that is rigorously independent of the reaction and is the definition used in practice~\cite{lecture9,transcript_l9}.

The same conservation argument shows when \emph{enthalpy} may be treated as equivalent to mixture fraction: in an adiabatic system with equal diffusivities of all species and heat (unity Lewis numbers), the absolute enthalpy is itself a conserved scalar with no source term and the same boundary values structure as $Z$, so it becomes a linear function of $Z$ --- the basis for parametrising the entire thermochemical state by $Z$ alone~\cite{lecture9,poinsot2011theoretical}. In a turbulent flame one solves not only for the mean $\favg Z$ but also for its \emph{variance} $\widetilde{Z^{\prime\prime 2}}$, because the variance fixes the width of the presumed PDF used to average the chemistry (Chapter~\ref{ch:turbcomb}); the production of variance by the mean gradient is balanced by its destruction, which is precisely the (mean) scalar dissipation rate $\widetilde\chi$ of Eq.~\eqref{eq:chi} --- so $\chi$ is the sink of mixture-fraction variance, linking micro-mixing directly to the spread of states seen in the scatter plots~\cite{lecture9,veynante2002turbulent}.
\section{The Burke--Schumann limit: ``mixed is burned''}

The simplest closed description assumes one-step, infinitely fast chemistry, so that fuel and oxidiser cannot coexist: wherever they meet they react completely and instantly --- the ``mixed is burned'' approximation~\cite{lecture9,poinsot2011theoretical}. With constant $c_p$ this gives the \textbf{Burke--Schumann} flame structure, in which the fuel and oxidiser mass fractions fall linearly to zero from their respective boundaries and the temperature rises linearly to a peak exactly at $Z=\Zst$~\cite{lecture9}. Whether the single step is treated as irreversible or reversible changes only the rounding of the peak; in both cases the mixture burns immediately on contact, and the entire thermochemical state becomes a function of $Z$ alone~\cite{lecture9}. A two-dimensional DNS of a reacting mixing layer confirms the picture: the temperature is highest along the stoichiometric mixture-fraction contour~\cite{lecture9}.

\begin{figure}[H]
\centering
\includegraphics[width=0.6\linewidth]{fig_mixfrac}
\caption{Mixture-fraction description of a non-premixed flame: $Z$ runs from 0 in the oxidiser to 1 in the fuel, and the flame sits at $Z=\Zst$~\cite{lecture9}.}
\label{fig:fig_mixfrac}
\end{figure}

\section{Departures from equilibrium and the flamelet model}

Real flames depart from the Burke--Schumann line because chemistry is finite and diffusion matters. Plotting a scalar (temperature or a species) against $Z$ for every point in a flame produces a \emph{scatter} that lies below the equilibrium line; the higher the local reaction rate relative to the strain, the closer the points cluster to equilibrium, and increasing the jet velocity or Reynolds number increases the scatter and the incidence of local extinction --- as seen in the canonical Sandia D/E/F flame measurements~\cite{lecture9}. To represent this simply we seek, as in the premixed case, a one-dimensional description of the structure normal to the flame.

This is achieved by the \textbf{counterflow flamelet model}~\cite{lecture9,peters2000turbulent}. One changes independent variables from physical coordinates $(x_1,x_2,x_3)$ to $(Z,x_\perp,x_\perp')$, with the first coordinate aligned with the direction of increasing $Z$ (normal to the flame) and the others parallel to it. Because $Z$ has no chemical source, the convective terms cancel in this transform (unlike the progress-variable case), and --- neglecting transport parallel to the flame, the \emph{flamelet assumption}, valid away from very high Karlovitz numbers --- the species and temperature equations reduce to a balance between reaction and a diffusion term governed by the \textbf{scalar dissipation rate}
\begin{equation}
\chi=2D\,\grad Z\cdot\grad Z=2D\left|\grad Z\right|^2 ,
\label{eq:chi}
\end{equation}
so that the flamelet equations read $\rho\,\dfrac{\chi}{2}\,\dfrac{\partial^2 Y_k}{\partial Z^2}+\dot\omega_k=0$ (and similarly for temperature)~\cite{lecture9}. The scalar dissipation rate characterises the rate of micro-mixing at a given $Z$ and is proportional to the imposed strain rate; an approximate profile $\chi(Z)$ from laminar counterflow studies peaks near the flame, with the key parameter being its stoichiometric value $\chi_{st}=\chi(\Zst)$~\cite{lecture9,poinsot2011theoretical}. Solving the flamelet equations for a given $\chi(Z)$ yields the profiles $Y_k(Z;\chi)$ and $T(Z;\chi)$, i.e.\ the flame structure parametrised by a single mixing parameter~\cite{lecture9}.

\paragraph{Limiting cases.} The flamelet solution has two instructive limits~\cite{lecture9,transcript_l9}. As $\chi\to0$ (low strain, slow mixing) the diffusion term vanishes and the solution approaches \emph{chemical equilibrium} --- the same as the Burke--Schumann / no-diffusion limit. As $\chi$ becomes very large (high strain, fast mixing) the reaction cannot keep pace with the diffusive losses, the flame becomes thinner and finally \emph{extinguishes}, leaving pure mixing without combustion. Increasing $\chi$ thus drives a diffusion flame toward extinction in exactly the way increasing stretch does for a premixed counterflow flame~\cite{lecture9}.

\section{Why diffusion flames respond immediately to strain}

The contrast with premixed flames is fundamental and worth stating explicitly~\cite{lecture9,transcript_l9}. A premixed flame has an intrinsic speed and can \emph{reposition} itself to escape an unfavourable strain, so its temperature and composition change only gradually with stretch and (absent differential diffusion) drop sharply only near extinction. A diffusion flame has no intrinsic speed: it is pinned to the $Z=\Zst$ surface set by the mixing field, so any change in strain immediately changes $\chi$ and therefore immediately changes the local structure --- there is no escape. At low to moderate strain, better mixing can initially \emph{raise} the temperature (more reactedness toward stoichiometric); at high strain the rising $\chi$ drives the flame toward extinction, and because the response is immediate the extinction strain rates are smaller in magnitude than for premixed flames~\cite{lecture9}. Differential diffusion still matters for diffusion flames, but it acts indirectly, by altering the mixing field $Z$ itself, rather than being directly coupled to the stretch as in the premixed case~\cite{lecture9,transcript_l9}.

\section{The candle: the canonical laminar diffusion flame}

The everyday laminar diffusion flame is the candle, and its anatomy illustrates the chapter~\cite{lecture5}. Heat from the flame melts the wax, which is drawn up the wick by capillary action and vaporised; the fuel vapour diffuses outward while ambient oxygen diffuses inward, and they burn where they meet in stoichiometric proportion. The flame is elongated into a teardrop by \emph{natural convection} --- hot gas rises and draws in cool air at the base --- so in microgravity, with oxygen supplied only by diffusion, a candle flame is nearly hemispherical~\cite{lecture5}. The visible structure shows the diffusion-flame physics: an inner dark zone where pyrolysis cracks the fuel into carbon particles under oxygen starvation, a bright yellow luminous zone where those soot particles glow (the yellow colour is incandescent soot, not chemiluminescence), and an outer veil where the soot finally oxidises in the oxygen-rich surroundings; if a draught cools the tip below about \SI{1000}{\celsius} the soot escapes unburnt~\cite{lecture5}. These features --- burning near $\phi\approx1$ at the flame surface, intrinsic stability, freedom from flashback, but a propensity to soot --- summarise why diffusion flames are safe and easy to control but tend to be high in soot and \ch{NO_x}, the trade-off that drives combustor design in Chapters~\ref{ch:emissions} and~\ref{ch:combustors}~\cite{lecture5}.
